\chapter{Radian et cercle trigonométrique}
\section*{Exercice 14.1 -- Arc et radian}
Un cercle a un rayon de $3$ cm. Quelle est la longueur d'un arc d'angle $2\pi/3$ ? Convertir cet angle en degrés.

\section*{Exercice 14.2 -- Valeur remarquable démontrée}
Retrouver $\cos(\pi/4)$ et $\sin(\pi/4)$ sans citer une table.

\section*{Exercice 14.3 -- Angles associés}
Calculer $\cos(5\pi/6)$, $\sin(5\pi/6)$, $\cos(-\pi/3)$ et $\sin(7\pi/6)$.

\section*{Exercice 14.4 -- Enroulement}
Donner un représentant dans $[0,2\pi[$ des angles $17\pi/3$ et $-9\pi/4$.

\section*{Exercice 14.5 -- Équation trigonométrique}
Résoudre $\cos x=1/2$ dans $\R$, puis dans $[0,2\pi[$.

\section*{Exercice 14.6 -- Archimède}
À partir de l'hexagone de côté $1$ inscrit au cercle unité, donner les bornes de $\pi$, puis le côté du dodécagone inscrit.
